\section{Number Theory}
\subsection{GCD Properties}
\begin{itemize}\setlength{\itemsep}{1pt}\setlength{\parskip}{0pt}\setlength{\parsep}{0pt}
  \item \textbf{Divisibility Trick:} If $a$ perfectly divides $(b \times c)$, and the GCD of $a$ and $b$ is $d$, then the remaining part of $a$ (which is $a/d$) MUST perfectly divide $c$.
  \item \textbf{Exponents:} The GCD of $(n^a - 1)$ and $(n^b - 1)$ is beautifully simple: it is just $n^{\gcd(a,b)} - 1$.
  \item \textbf{Geometry:} If you draw a straight line on a 2D grid from $(0,0)$ to $(a,b)$, the number of exact grid points strictly between them is $\gcd(|a|, |b|) - 1$.
\end{itemize}

\subsection{Extended Euclid, CRT, Diophantine}
\texttt{extgcd}: $ax+by=g=\gcd(a,b)\ge 0$. $Ax+By=C$: $x=x_0+k\frac{B}{g}, y=y_0-k\frac{A}{g}$ ($g\mid C$). \textbf{Frobenius:} $\gcd(a,b)=1\implies$ max unreachable $ab-a-b$; count $\frac{(a-1)(b-1)}{2}$. \textbf{CRT:} general moduli.
\begin{minted}{c++}
namespace LDE {
  using i128 = __int128_t;
  // in: a, b, &x, &y | out: g=gcd>=0, a*x+b*y=g
  ll extgcd(ll a, ll b, ll &x, ll &y) {
    if (!b) return x = (a < 0 ? -1 : 1), y = 0, abs(a);
    ll g = extgcd(b, a % b, y, x);
    return y -= a / b * x, g;
  }
  // in: a, m | out: inv in [0,m-1] s.t. (a*x)%m=1
  // returns -1 if gcd(a,m) != 1 or m <= 1
  ll mod_inv(ll a, ll m) {
    if (m <= 1) return -1;
    ll x, y, g = extgcd((a % m + m) % m, m, x, y);
    return g == 1 ? (x % m + m) % m : -1;
  }
  // Solves: a*x = b (mod m)
  // in: a, b, m
  // out: {x0, step} => x = x0 (mod step).
  // Sols in [0,m-1]: x0 + k*step (total m/step).
  // Returns {-1, -1} if no sol (b % gcd != 0)
  pair<ll, ll> congruence(ll a, ll b, ll m) {
    if (m <= 0) return {-1, -1};
    ll u, v, g = extgcd(a, m, u, v);
    if (b % g) return {-1, -1};
    ll st = m / g, x0 = ((i128)u * (b / g) % st + st) % st;
    return {x0, st};
  }
  // Solves: x = r1 (mod m1), x = r2 (mod m2)
  // in: r1, m1, r2, m2 (any m1, m2 > 0)
  // out: {x0, lcm} => x = x0 (mod lcm), 0<=x0<lcm
  // Returns {-1, -1} if inconsistent
  pair<ll, ll> crt(ll r1, ll m1, ll r2, ll m2) {
    if (m1 <= 0 || m2 <= 0) return {-1, -1};
    r1 = (r1 % m1 + m1) % m1; r2 = (r2 % m2 + m2) % m2;
    auto [k0, st] = congruence(m1, r2 - r1, m2);
    if (k0 == -1) return {-1, -1};
    i128 lcm = (i128)m1 * st;
    return {(r1 + (i128)k0 * m1) % lcm, (ll)lcm};
  }
  // in: vector r, m | out: {x0, lcm} / {-1,-1}
  pair<ll, ll> crt(const vector<ll>& r, const vector<ll>& m) {
    pair<ll, ll> c = {0, 1};
    for (int i = 0; i < (int)r.size(); i++)
      if ((c = crt(c.first, c.second, r[i], m[i])).second == -1) break;
    return c;
  }
  // Solves: a*x + b*y = c (Linear Diophantine)
  // in: a, b, c, &x, &y, &g
  // out: true if sol exists (g|c), sets g, (x,y)
  // Base sol: x in [0, |b/g|-1] (never overflows)
  // General: x' = x+k*(b/g), y' = y-k*(a/g)
  bool any_sol(ll a, ll b, ll c, ll &x, ll &y, ll &g) {
    if (!a && !b) { g = x = y = 0; return c == 0; }
    g = extgcd(a, b, x, y);
    if (c % g) return false;
    if (b) {
      ll m = abs(b / g);
      x = ((i128)x * (c / g) % m + m) % m;
      y = (c - (i128)a * x) / b;
    } else x = c / a, y = 0;
    return true;
  }
  i128 flr(i128 a, i128 b) { return a >= 0 ? a / b : (a - b + 1) / b; }
  i128 cil(i128 a, i128 b) { return a >= 0 ? (a + b - 1) / b : a / b; }
  pair<i128, i128> range_k(i128 u, i128 v, i128 L, i128 R) {
    if (v < 0) u = -u, v = -v, swap(L, R), L = -L, R = -R;
    return {cil(L - u, v), flr(R - u, v)};
  }
  // in: a*x + b*y = c and box [lx,rx] x [ly,ry]
  // out: count of integer pairs (x,y) in box
  // min x sol: k = (b/g > 0 ? kl : kr)
  ll count_sol(ll a, ll b, ll c, ll lx, ll rx, ll ly, ll ry) {
    if (lx > rx || ly > ry) return 0;
    if (!a && !b) return c ? 0 : (ll)((i128)(rx - lx + 1) * (ry - ly + 1));
    if (!a) return (c % b == 0 && ly <= c / b && c / b <= ry) ? rx - lx + 1 : 0;
    if (!b) return (c % a == 0 && lx <= c / a && c / a <= rx) ? ry - ly + 1 : 0;
    ll x, y, g;
    if (!any_sol(a, b, c, x, y, g)) return 0;
    auto [k1, k2] = range_k(x, b / g, lx, rx);
    auto [k3, k4] = range_k(y, -a / g, ly, ry);
    i128 kl = max(k1, k3), kr = min(k2, k4);
    return kl > kr ? 0 : (ll)(kr - kl + 1);
  }
}
\end{minted}
\subsection{SPF, Möbius, Primes}
\begin{minted}{c++}
// O(N)
const int N = 1e6 + 9;
int spf[N];      // Smallest prime factor
int phi[N];      // Euler's totient function
int mob[N];      // Möbius function
int divisor[N];  // Number of divisors d(n)
int cnt[N];      // Exponent of spf in n
vector<int> primes;

void linear_sieve() {
  phi[1] = mob[1] = divisor[1] = 1;
  for (int i = 2; i < N; i++) {
    if (!spf[i]) {
      spf[i] = i; primes.push_back(i);
      phi[i] = i - 1; mob[i] = -1;
      divisor[i] = 2; cnt[i] = 1;
    }
    for (int p : primes) {
      if (1LL * i * p >= N) break;
      spf[i * p] = p;
      if (i % p == 0) {
        phi[i * p] = phi[i] * p; mob[i * p] = 0;
        cnt[i * p] = cnt[i] + 1;
        divisor[i * p] = divisor[i] / (cnt[i] + 1)
                         * (cnt[i * p] + 1);
        break;
      } else {
        phi[i * p] = phi[i] * (p - 1);
        mob[i * p] = -mob[i]; cnt[i * p] = 1;
        divisor[i * p] = divisor[i] * 2;
      }
    }
  }
}
// factorize x (x <= (N-1)^2)
vector<pair<int,int>> factorize(int x) {
  vector<pair<int,int>> ret;
  for (int p : primes) {
    if (1LL * p * p > x) break;
    if (x % p == 0) { int c = 0; while (x % p == 0) x /= p, c++; ret.emplace_back(p, c); }
  }
  if (x > 1) {
    if (x >= N) ret.emplace_back(x, 1);
    else { while (x > 1) { int p = spf[x], c = 0; while (x % p == 0) x /= p, c++; ret.emplace_back(p, c); } }
  }
  return ret;
}
// divisors of x (x <= (N-1)^2)
vector<int> divisors(int x) {
  vector<int> ret = {1};
  for (auto [p, c] : factorize(x)) {
    int n = ret.size(); ll a = 1;
    for (int d = 1; d <= c; d++) { a *= p; for (int i = 0; i < n; i++) ret.push_back(ret[i] * a); }
  }
  return ret;
}
\end{minted}
\subsection{SOD \& POD (Divisor Functions)}
\begin{itemize}\setlength{\itemsep}{1pt}\setlength{\parskip}{0pt}\setlength{\parsep}{0pt}
  \item \textbf{Divisor Counts:} For $n = \prod p_i^{e_i}$, positive divisor count $d(n) = \prod (e_i+1)$. \\
  $\omega(n) =$ distinct prime factors count; $\Omega(n) = \sum e_i$ (with multiplicity).
  \item \textbf{SOD (Sum of Divisors):} $\sigma(n) = \sum_{d \mid n} d = \prod \frac{p_i^{e_i+1}-1}{p_i-1} = \prod (1+p_i+\dots+p_i^{e_i})$.
  \item \textbf{$k$-th Power Divisor Sum:} $\sigma_k(n) = \sum d^k = \prod \frac{p_i^{k(e_i+1)}-1}{p_i^k-1}$.
  \item \textbf{Proper Divisor Sum:} $s(n) = \sigma(n) - n$: $=n$ (perfect), $>n$ (abundant), $<n$ (deficient).
  \item \textbf{POD (Product of Divisors):} $\text{POD}(n) = \prod_{d \mid n} d = n^{d(n)/2} = (\sqrt{n})^{d(n)}$. \\
  Each pair $(d, n/d)$ multiplies to $n$. When $d(n)$ is odd ($n$ is square), use $(\sqrt{n})^{d(n)}$ or $\prod p_i^{e_i \cdot d(n)/2}$ to avoid float error. Ex: $12 \to d=6, \sigma=28, \text{POD}=12^3$.
  \item \textbf{Mod SOD without Division:} Compute $S(k) = \sum_{j=0}^{k-1} r^j \pmod M$ when $\gcd(r-1, M) \ne 1$ (no mod inverse) via D\&C in $O(\log k)$: $S(2m) = S(m)(1+r^m)$, $S(2m+1) = S(2m) + r^{2m}$ (base $S(0)=0, S(1)=1$).
  \item \textbf{Legendre's Formula:} $v_p(n!) = \sum_{k \ge 1} \lfloor n/p^k \rfloor$ is the exponent of prime $p$ in $n!$. \\
  Trailing zeros of $n!$: base 10 is $v_5(n!)$; general base $B = \prod q_j^{a_j}$ is $\min_j \lfloor v_{q_j}(n!)/a_j \rfloor$.
  \item \textbf{Wilson's Theorem:} $n > 1$ is prime $\iff (n-1)! \equiv -1 \pmod n$. \\
  Compute $(p-k)! \pmod p$ for small $k \ll p$: $(p-k)! \equiv \frac{(-1)^k}{(k-1)!} \pmod p$ (e.g. $(p-2)! \equiv 1$).
  \item \textbf{Fermat \& Euler's Criterion (QR):} If $\gcd(a,p)=1$, $a^{p-1} \equiv 1 \pmod p \implies a^{-1} \equiv a^{p-2}$. \\
  $a^{(p-1)/2} \equiv \left(\frac{a}{p}\right) \pmod p$: $+1 \implies a$ is quadratic residue (square mod $p$), $-1 \implies$ non-residue (no square root mod $p$), $0 \implies p \mid a$.
  \item \textbf{Carmichael $\lambda(n)$:} Smallest exponent with $a^{\lambda(n)} \equiv 1 \pmod n$ for all $\gcd(a,n)=1$ (divides $\varphi(n)$). For odd prime powers: $\lambda(p^k) = p^{k-1}(p-1) = \varphi(p^k)$. For powers of 2: $\lambda(2)=1, \lambda(4)=2, \lambda(2^k)=2^{k-2}$ ($k \ge 3$). General: $\lambda(n) = \operatorname{lcm}_i \lambda(p_i^{e_i})$. Ex: $\lambda(1000) = \operatorname{lcm}(2, 100) = 100 < \varphi(1000)=400$.
  \item \textbf{Primitive Root:} Generator $g$ whose powers produce all $\varphi(n)$ coprime residues mod $n$ (order $\varphi(n)$). Exists iff $n \in \{1, 2, 4, p^k, 2p^k\}$ (odd prime $p$); count is $\varphi(\varphi(n))$. \\
  Test: $g$ is primitive iff $g^{\varphi(n)/q} \not\equiv 1 \pmod n$ for all prime factors $q \mid \varphi(n)$.
\end{itemize}
\subsection{Segmented Sieve}
Key idea: Sieve small primes $\le \sqrt{R}$, then strike composites in $[L,R]$ using offsets. Efficient for very large ranges ($L,R \le 10^{12}$, $R-L \le 10^7$).
\begin{minted}{c++}
vector<char> segmentedSieve(ll L, ll R) {
  ll lim = sqrt(R);
  vector<char> mark(lim + 1, false);
  vector<ll> primes;
  for (ll i = 2; i <= lim; ++i) {
    if (!mark[i]) {
      primes.emplace_back(i);
      for (ll j = i * i; j <= lim; j += i) mark[j] = true;
    }
  }
  vector<char> isPrime(R - L + 1, true);
  for (ll i : primes) {
    ll sm = (L / i) * i;
    if (sm < L) sm += i;
    for (ll j = max(i * i, sm); j <= R; j += i)
      isPrime[j - L] = false;
  }
  if (L <= 1)
    for (ll i = L; i <= 1 && i <= R; i++) isPrime[i - L] = false;
  return isPrime;
}
\end{minted}
\subsection{Miller-Rabin Primality Test}
Deterministic for all 64-bit integers. $O(\log^3 n)$. Uses fixed bases $\{2, 325, 9375, 28178, 450775, 9780504, 1795265022\}$. Returns \texttt{true} if $n$ is prime, \texttt{false} if composite.
\begin{minted}{c++}
ull pow(ull a, ull t, ull mod) {
  ull r = 1;
  for (; t; t >>= 1, a = (__uint128_t)a * a % mod)
    if (t & 1) r = (__uint128_t)r * a % mod;
  return r;
}
bool isprime(ull n) {
  if (n < 2 || n % 6 % 4 != 1) return (n | 1) == 3;
  ull A[] = {2, 325, 9375, 28178, 450775,
  9780504, 1795265022},
  s = __builtin_ctzll(n - 1), d = n >> s;
  for (ull a : A) {
    ull p = pow(a % n, d, n), i = s;
    while (p != 1 && p != n - 1 && a % n && i--)
      p = (__uint128_t)p * p % n;
    if (p != n - 1 && i != s) return 0;
  }
  return 1;
}
\end{minted}
\subsection{Pollard Rho (Fast Factor)}
Factor any 64-bit integer in $O(n^{1/4})$ expected time. Uses Montgomery multiplication for fast modular arithmetic. \texttt{find\_factor(n)} returns a non-trivial factor (or $n$ if prime). Pair with Miller-Rabin: if \texttt{isprime(n)} $\to$ done, else factor.
\textbf{Note:} Randomized. Can get stuck in an infinite loop if $N$ is prime, always check \texttt{isprime(N)} first.
\begin{minted}{c++}
#define ull uint64_t
struct Montgomery {
  ull n, nr;
  Montgomery(ull n) {
    this->n = n, nr = 1;
    for (int i = 0; i < 6; i++)
      nr *= 2 - n * nr;
  }
  ull reduce(__uint128_t x) const {
    ull q = ull(x) * nr,
    m = ((__uint128_t)q * n) >> 64;
    ull res = (x >> 64) + n - m;
    if (res >= n) res -= n;
    return res;
  }
  ull multiply(ull x, ull y) {
    return reduce((__uint128_t)x * y); }
};
ull f(ull x, ull n, ull a) {
  return ((__uint128_t)x * x + a) % n; }
ull diff(ull a, ull b) {
  return a > b ? a - b : b - a; }
const int M = 1024;
const ull SEED = 42;
ull find_factor(ull n, ull x0 = 2, ull a = 1) {
  ull x = x0, y = x0, d = 1;
  while (d == 1) {
    x = f(x, n, a);
    y = f(f(y, n, a), n, a);
    d = gcd(diff(x, y), n);
  }
  return d == n ? find_factor(n, x0, a + 1) : d;
}
\end{minted}
\subsection{Phi (Euler Totient)}
$\varphi(n)$ = count of integers in $[1,n]$ coprime to $n$. Formula: $\varphi(n) = n \prod_{p|n} (1-1/p)$. Computed via trial division. \texttt{phi\_1\_to\_n()} computes $\varphi$ for all $i \in [1,N]$ using linear sieve: $O(N)$.
\begin{minted}{c++}
int phi(int n) {
  int result = n;
  for (int i = 2; i * i <= n; i++) {
    if (n % i == 0) {
      while (n % i == 0) n /= i;
      result -= result / i;
    }
  }
  if (n > 1) result -= result / n;
  return result;
}
void phi_1_to_n() {
  phi[1] = 1;
  for (int i = 2; i < N; i++) {
    if (phi[i] == 0) {
      primes.push_back(i);
      phi[i] = i - 1;
    }
    for (int p : primes) {
      if (i * p >= N) break;
      if (i % p == 0) {
        phi[i * p] = phi[i] * p; break;
      } else {
        phi[i * p] = phi[i] * (p - 1);
      }
    }
  }
}
\end{minted}
Sum of numbers $\le x$ coprime to $x$: $\displaystyle\sum_{\substack{1 \le i \le x \\ \gcd(i, x) = 1}} i = \frac{x \cdot \phi(x)}{2}$.

\textbf{Euler's Totient Theorem}: If $\gcd(a,n)=1$, $a^{\phi(n)} \equiv 1 \pmod n$, and consequently $a^n \equiv a^{n-\phi(n)} \pmod n$. For general $a,n$ (not necessarily coprime), for $m \ge \log_2 n$: $a^m \equiv a^{(m \bmod \phi(n)) + \phi(n)} \pmod n$.

\subsection{Möbius Function and Inversion}
The Möbius function $\mu(n)$: \\
\textbullet{} $\mu(n) = 1$ if $n$ is square-free with even number of prime factors. \\
\textbullet{} $\mu(n) = -1$ if $n$ is square-free with odd number of prime factors. \\
\textbullet{} $\mu(n) = 0$ if $n$ has a squared prime factor.

\textbf{Properties \& Inversion}: \\
\textbullet{} $\sum_{d|n} \mu(d) = [n = 1]$ \\
\textbullet{} $g(n) = \sum_{d|n} f(d) \iff f(n) = \sum_{d|n} \mu(d)g(n/d)$ \\
\textbullet{} $\phi(n) = \sum_{d|n} \mu(d) \frac{n}{d}$ and $n = \sum_{d|n} d \cdot \mu\left(\frac{n}{d}\right)$ \\
\textbullet{} $\frac{n}{\phi(n)} = \sum_{d|n} \frac{\mu(d)}{d}$

\textbf{GCD Sums}: \\
\textbullet{} $\sum_{i=1}^n [\gcd(i, n) = k] = \phi\left(\frac{n}{k}\right)$ \\
\textbullet{} $\sum_{i=1}^n \gcd(i, n) = \sum_{d|n} d \cdot \phi\left(\frac{n}{d}\right)$ \\
\textbullet{} $\sum_{i=1}^n \frac{1}{\gcd(i,n)} = \frac{1}{n} \sum_{d|n} d \cdot \phi(d)$ \\
\textbullet{} $\sum_{i=1}^n \frac{i}{\gcd(i,n)} = \frac{1}{2} \sum_{d|n} d \cdot \phi(d)$

\subsection{Gauss Circle Theorem}
Lattice points inside $m^2 + n^2 \le r^2$: \\
$N(r) = 1 + 4 \sum_{i=0}^\infty \left( \lfloor \frac{r^2}{4i+1} \rfloor - \lfloor \frac{r^2}{4i+3} \rfloor \right)$ \\
\textbf{Note:} Counts integer points $(x,y)$ such that $x^2 + y^2 \le r^2$.

\subsection{Sum of Floor($n/i$)}
$\displaystyle\sum_{i=1}^n \left\lfloor\frac{n}{i}\right\rfloor$ in $O(\sqrt{n})$ by grouping equal terms.
\textbf{Applications:} Divisor sum $\sum d(i)$, Dirichlet hyperbola basis, $\sum_{i=1}^n \lfloor n/i\rfloor \phi(i)$.
\begin{minted}{c++}
ll sum_floor(ll n) {
  ll ans = 0;
  for (ll l = 1, r; l <= n; l = r + 1) {
    r = n / (n / l);
    ans += (r - l + 1) * (n / l);
  }
  return ans;
}
\end{minted}

\subsection{Floor Sum}
Computes $\displaystyle\sum_{i=0}^{n-1} \left\lfloor\frac{a\cdot i + b}{m}\right\rfloor$ in $O(\log m)$ via Euclidean recursion. From AtCoder Library (ACL).
\textbf{Applications:} Lattice counting, $\sum (a\cdot i + b) \% m$, digit DP.
\begin{minted}{c++}
ll floor_sum(ll n, ll m, ll a, ll b) {
  ll ans = 0;
  if (a >= m) {
    ans += (n - 1) * n * (a / m) / 2;
    a %= m;
  }
  if (b >= m) {
    ans += n * (b / m);
    b %= m;
  }
  ll y_max = (a * n + b) / m;
  if (y_max == 0) return ans;
  ans += floor_sum(y_max, a, m, (a * n + b) % m);
  return ans;
}
\end{minted}

\subsection{Discrete Logarithm (Baby-Step Giant-Step)}
Solves $a^x \equiv b \pmod m$ for smallest $x \ge 0$. Requires $\gcd(a,m)=1$. $O(\sqrt{m})$ meet-in-the-middle.
\begin{minted}{c++}
ll discrete_log(ll a, ll b, ll m) {
  a %= m; b %= m;
  if (b == 1) return 0;
  ll n = sqrt(m) + 1;
  unordered_map<ll, ll> baby;
  ll cur = b;
  for (ll q = 0; q < n; q++) {
    baby[cur] = q;  // overwrite: keep the largest q for smallest x
    cur = cur * a % m;
  }
  ll an = 1;
  for (ll i = 0; i < n; i++) an = an * a % m;
  cur = 1;
  for (ll p = 0; p <= n; p++) {
    if (baby.count(cur)) {
      ll x = p * n - baby[cur];
      if (x >= 0) return x;
    }
    cur = cur * an % m;
  }
  return -1;
}
\end{minted}

\subsection{Lifting The Exponent (LTE)}
\textbf{What it is:} Finds the highest power of a prime $p$ dividing $a^n \pm b^n$. $\nu_p(X)$ is the exponent of $p$ in $X$'s prime factorization. 
\textbf{When to use:} Checking if $p^k \mid a^n \pm b^n$ for huge exponents, or finding trailing zeros.
\textbf{How to use:} Check strict preconditions ($p \mid a \pm b$, $p \nmid a$, $p \nmid b$), then add powers:

\textbf{Odd Prime ($p \ge 3$):}
\begin{itemize}
  \item $\nu_p(a^n - b^n) = \nu_p(a-b) + \nu_p(n)$
  \item $\nu_p(a^n + b^n) = \nu_p(a+b) + \nu_p(n)$ \quad (requires $n$ to be \textbf{odd})
\end{itemize}

\textbf{Even Prime ($p = 2$), requires $a,b$ \textbf{odd}:}
\begin{itemize}
  \item $\nu_2(a^n - b^n) = \nu_2(a-b) + \nu_2(a+b) + \nu_2(n) - 1$ \quad (if $n$ is \textbf{even})
  \item $\nu_2(a^n - b^n) = \nu_2(a-b)$ \quad (if $n$ is \textbf{odd})
  \item $\nu_2(a^n + b^n) = 1$ \quad (if $n$ is \textbf{even})
  \item $\nu_2(a^n + b^n) = \nu_2(a+b)$ \quad (if $n$ is \textbf{odd})
\end{itemize}

\subsection{Fibonacci Number Properties}
$F_0 = 0, F_1 = 1, F_n = F_{n-1} + F_{n-2}$.
\begin{itemize}
  \item To check if one Fibonacci divides another, just check their indices: $F_m \mid F_n \iff m \mid n$.
  \item Calculate the GCD of their indices instead: $\gcd(F_m, F_n) = F_{\gcd(m,n)}$.
  \item Simplifies alternating sum puzzles or determinants: $F_{n-1}F_{n+1} - F_n^2 = (-1)^n$.
  \item Backbone for $O(\log N)$ doubling/matrix algorithms combining states $n,m$: $F_{n+m} = F_{n-1}F_m + F_nF_{m+1}$.
  \item Maps recursive steps that go below zero back to positive: $F_{-n} = (-1)^{n+1}F_n$.
  \item Eliminates $O(N)$ loops when querying the sum of $N$ terms: $\sum_{i=1}^n F_i = F_{n+2} - 1$.
  \item Eliminates $O(N)$ loops in grid/tiling sums: $\sum_{i=1}^n F_i^2 = F_nF_{n+1}$.
  \item For huge $N$ and small modulo $m$, the sequence repeats. The cycle length is $\le 6m$. Compute $N \bmod \text{cycle}$.
  \item $O(1)$ direct calc. Usable under modulo only if $5$ is a quadratic residue mod $p$: $F_n = \frac{\phi^n - \psi^n}{\sqrt{5}}$ ($\phi = \frac{1+\sqrt{5}}{2},\ \psi = \frac{1-\sqrt{5}}{2}$).
\end{itemize}

\begin{minted}{c++}
// Fast Doubling — returns {F_n, F_{n+1}} in O(log n),  n up to 1e18
// F_0=0, F_1=1.  F_n = fib(n).first
const ll MOD = 1e9 + 7;
pair<ll, ll> fib(ll n) {
  if (n == 0) return {0, 1};
  auto [a, b] = fib(n >> 1);
  ll c = a * ((2 * b % MOD - a + MOD) % MOD) % MOD;
  ll d = (a * a % MOD + b * b % MOD) % MOD;
  if (n & 1) return {d, (c + d) % MOD};
  return {c, d};
}
\end{minted}
